\section{\label{II_1_A}Traitement audiovisuel automatique}


%marquer davantage la vision par ordinateur classique et la vision par ordinateur fondée sur le deep learning, bcp utilisée dans l'analyse de l'audiovisuel

\subsection{État des lieux}

Les technologies appliquées à l'analyse de contenus vidéo et audio, comprenant la computer vision et la reconnaissance de discours automatique (ASR) sont globalement désignées sous le terme de \gls{AudiovisualProcessing} (AVP)\footcite{noord2021}.

Appliquées à un grand nombre de données, ces technologies permettent de passer du \textit{"close"} au \textit{"distant viewing"}, c’est-à-dire d’une vue des données rapprochée, réduite, à une vue d’ensemble. \footcite{noord2021} Les tâches de traitement automatique permettraient en effet de passer de l’analyse manuelle de certaines archives audiovisuelles, à une vue d'ensemble sur les données afin d'en extraire les tendances globales, modifiant ainsi notre rapport à ces documents. 

Les tâches d’analyse automatique pouvant être appliquées aux documents audiovisuels sont multiples, regroupant notamment le traitement du signal sonore (reconnaissance des locuteurs, transcription et classification audio) et l’analyse d’images (reconnaissance faciale, détection et classification d’objets, extraction de textes). Les chercheurs analysant les objets visuels s’orientent de plus en plus vers le contenu audiovisuel, puisque la détection du mouvement constitue un levier déterminant pour la mise en place d’autres tâches automatiques de reconnaissance \footcite{brostow2009}. 

Toutefois, l’accès au contenu actionné par le flux vidéo exige une structuration spécifique. Si la détection automatique des plans dans un flux vidéo est maîtrisée puisqu’elle n’engage que des changements visibles techniquement dans l’image, la détection sémantique de séquences reste inexploitée. Comme le souligne Samarth Bhargav, “pouvoir détecter de manière fiable les limites des scènes constituerait un exploit technique incroyable, qui ouvrirait d'énormes perspectives pour l'analyse sémantique et narrative des séquences vidéo et cinématographiques. [Traduction personnelle]" \footnote{\textit{"reliably being able to detect scene boundaries would be an incredible technical feat, which would create tremendous potential for semantic and narrative analysis of video and film material"} \cite{bhargav2019}}.

En effet, de nombreux projets sont dédiés à l’analyse automatique de corpus audiovisuels, toutefois, aucun projet n’est concrètement dédié à la \gls{segtemp} en différentes \gls{sequence} afin d'en indexer le contenu. L’Institut National Audiovisuel (\gls{ina}), est notamment pionnier dans la proposition de projets d’analyse automatique appliquée à l’audiovisuel. L’\gls{ina} a par exemple mis en place un outil de segmentation, \textit{InaSpeechSegmenter}, toutefois dédié à la segmentation du discours entre hommes et femmes, dans une optique statistique \footcite{zotero-item-787}. Cet algorithme de segmentation, appliqué au son, ne peut être pris comme exemple d’application pour notre étude, centrée sur l’analyse du sens de l’image en mouvement.  

Par ailleurs, le projet \textit{MeMad}, réunit un ensemble d’outils dédiés à l’analyse automatique de l’audiovisuel développés par l’Université d’Aalto, EURECOM, l’INA et l'Université d'Helsinki \footcite{zotero-item-790}. Parmi les outils développés, nous pouvons citer par exemple un outil de traduction du discours (\textit{Speech translation}), l’outil \textit{InaFaceGender}, un outil de reconnaissance faciale (\textit{Face recognition}), un outil permetttant d’identifier la langue parlée (\textit{Spoken language identification}) \footcite{zotero-item-790}. Ces projets restent centrés sur l’analyse du son et du texte inclut dans l’image et ne peuvent donc pas être pris en compte dans notre analyse. 
Au sein de cette sous section, nous tenterons ainsi d'évaluer les outils et méthodes disponibles et dédiés au traitement audiovisuels, afin de proposer un processus adéquat. 

%https://www.ai4media.eu/  

%Groupe de recherche et de réflexion et de développement d’outils pour l’analyse des médias par l’IA : pas de modules dédiés à la segmentation par séquence. Modules développés : reconnaissance d’objets, génération de résumés, segmentation sémantique. Mais modules bcp basés sur les mécanismes d’auto-attention pour ‘l'analyse vidéo. Documentation des projets non accessible. Un outil intéressant de détection de chnagement de plans mais documentation non accessible : 
%source: voir word 

\subsection{La \gls{Vision par ordinateur}}

\subsubsection{La \gls{Vision par ordinateur} classique}

L’application des technologies d’intelligence artificielle au domaine des images (images fixe ou vidéo) est désignée sous le terme de \textit{computer vision}, \gls{Vision par ordinateur} en français. Le terme définit un sous-domaine d’application de l’intelligence artificielle qui permet d’entraîner un système d’\gls{ia} à l’analyse des images fixes ou en mouvement \footcite{stryker2025}. 

La \gls{Vision par ordinateur} s’appuie largement sur des modèles d' \gls{Apprentissage machine} (\textit{machine learning}), qui désigne une branche de l'\gls{ia} regroupant les techniques, méthodes et processus nécessaires afin de rendre une machine ou un système informatique plus intelligent \footcite{zotero-item-584}. L'\gls{Apprentissage machine} permet d’entraîner des algorithmes capables d’apprendre à partir de données d’entraînement, afin de prédire ou de prendre des décisions sur de nouvelles données \footcite{zotero-item-702}. La \gls{Vision par ordinateur} classique s’appuie sur l’\gls{Apprentissage machine} précisément pour sa capacité à entraîner des systèmes à prédire. En effet, les corpus d’images fixes et d’images en mouvement sont des phénomènes complexes à appréhender pour des algorithmes, contrairement à des données statistiques et mathématiques. La nature de ces corpus les rend imprévisibles autant dans leur forme que dans leur contenu. Les algorithmes d’intelligence artificielle ne peuvent donc pas seulement réaliser une analyse prédictive de ce type de corpus, définie par des critères mathématiques \footcite{mahony2020}. 

Les corpus audiovisuels nécessitent une analyse prédictive, que les algorithmes sont entraînés à réaliser via des méthodes d'\gls{Apprentissage machine} \footcite{mahony2020}. De plus, la plupart de ces corpus sont d’une taille conséquente, forçant ainsi la mise en place d'algorithmes de prédiction, puisque l’entraînement machine ne peut être mis en place sur l’ensemble des données, seulement sur un échantillon donné. 

\subsubsection{\gls{Vision par ordinateur} fondée sur l'\gls{Apprentissage profond}}

Pour des tâches telles que la segmentation en différentes séquences, qui implique l’analyse du contenu de l’image, sont utilisés des modèles d’intelligence artificielle reposant sur des réseaux de neurones, et notamment des \gls{Réseaux de neurones profonds} \footcite{norindr2023}. Ces modèles relèvent également de la \gls{Vision par ordinateur}, mais fondée sur l'\gls{Apprentissage profond}, permettant notamment de gérer l'analyse de flux vidéo. La méthode de mise en place de ces \gls{Réseaux de neurones profonds} en plusieurs couches, est désignée sous le terme de \textit{deep learning}, \gls{Apprentissage profond}, une sous-catégorie de l'\gls{Apprentissage machine}. L'\gls{Apprentissage profond} est basé sur des réseaux de neurones artificiels inspirés du fonctionnement du cerveau humain. Les neurones sont en réalité des cellules de calcul organisées en différentes couches qui interagissent les unes avec les autres afin de s’adapter aux données reçues, et de prendre une décision quelconque. Davantage entrainé que programmé, un modèle fondé sur l’\gls{Apprentissage profond} permet d’exécuter des tâches automatisées, réduisant le besoin d'intervention humaine, et améliorant ainsi le rapport cout-efficacité \footcite{mahony2020}. Il suffirait de fournir au modèle quelques images annotées, dont le réseau de neurones repèrerait les caractéristiques à analyser sur de nouvelles données \footcite{mahony2020}. Par rapport aux techniques traditionnelles de \gls{Vision par ordinateur}, l'\gls{Apprentissage profond} permet d’obtenir une plus grande précision dans les tâches d'analyse audiovisuelle telles que la classification d'images, la segmentation, la détection d'objets \footcite{mahony2020}.

\subsection{\gls{Transformers} et \gls{cnn}}

Au sein des réseaux de neurones, le type de réseau spécifiquement employé dans le domaine de la \gls{Vision par ordinateur} pour la classification d’images et la reconnaissance d’objets, est celui des \gls{cnn}. Appliqués à des corpus d’images, les \gls{cnn} peuvent par exemple analyser une image selon une hiérarchie de caractéristiques afin d’en extraire les couleurs et les contours \footcite{brachmann2017}. Les tâches que permettent les \gls{cnn} sont plutôt bien adaptées à l’image fixe, puisqu’ils extraient des éléments fixes de l’image.  

Néanmoins, d’autres modèles issus de l'\gls{Apprentissage profond} tels que les \gls{Transformers} s’avèrent davantage efficace sur des tâches dédiées à l’image en mouvement, dont fait partie la segmentation. Contrairement aux \gls{cnn}, les \gls{Transformers} fonctionnent grâce à un mécanisme d’auto-attention qui leur permet de détecter les relations et dépendances entre les données \footcite{bergmann2025}. Tandis que les \gls{cnn} analysent les données de manière locale dans l'image, sans prendre en compte les interdépendances qui les régissent temporellement, les \gls{Transformers} examinent simultanément les données dans leur ensemble et décident de se concentrer sur certaines étapes de chacune d'entre elles \footcite{bergmann2025}. Les \gls{cnn} ne permettent pas, par exemple, de discerner les “dépendances à longue portée” telles que les corrélations entre les pixels entre différentes images qui ne seraient pas voisins \footcite{bergmann2025}. Ainsi, ils seraient davantage adaptés à l’analyse d’images en mouvement puisqu’ils prendraient en compte les relations entre les images, tandis que les \gls{cnn} seraient davantage adaptés aux images fixes.  Idéalement, la mise en place d’une segmentation automatique des archives audiovisuelles doit donc être fondée sur une architecture de type \gls{Transformers}.  


Toutefois, dans un contexte patrimonial, avec des ressources limitées, ces algorithmes fondés sur l’intelligence artificielle possèdent leurs limites. Premièrement, les réseaux de neurones sont conçus comme des boites noires, ne permettant pas d’interpréter les analyses qui sont produites \footcite{brachmann2017}. En effet, une fois entraînés, il est difficile d’en modifier les paramètres puisqu’ils sont entremêlés entre les différentes cellules de calcul. La réflexion produite est ainsi difficile à comprendre et à exposer\footcite{mahony2020}. 

Deuxièmement, les algorithmes de type \gls{Transformers} sont complexes et restent difficiles à adapter à un corpus audiovisuel lui-même complexe \footcite{yuan2025}. Ils demandent également de grandes puissances de calcul et un temps de paramétrage important, ce qui les rend inadaptés à des contextes de gestion de collections patrimoniales. De plus, ils requièrent un grand nombre de paramètres issu d’un large jeu de données d’entraînement afin d’éviter un surapprentissage impliquant des biais multiples \footcite{pattichis2023}. Ces modèles sont, de surcroit, généralement entraînés sur des jeux de donénes inadaptés, qu'il s'agirait de réadapter à un contexte patrimonial \footcite{pattichis2023}. Cette réadaptation impliquerait un temps conséquent dédié à la validation et à la correction de l'analyse automatique produite. Certaines approches algorithmiques non neuronales ou hybrides, qui combineraient \gls{Vision par ordinateur} et \gls{cnn} par exemple, seraient ainsi davantage adaptées car fondées sur des calculs et des seuils mathématiques les rendant plus stables. 
Ainsi, la nature de chaque modèle conditionne la qualité d'analyse automatique d'un flux vidéo. Afin de synthétiser les apports et les contraintes de ces différentes approches, le Tableau \ref{tab:comparatif_avp} ci-dessous permet de comparer les compétences des types de modèles cités pour l'analyse d'un flux vidéo. 


\begin{table}[H]
	\centering
	\small
	\renewcommand{\arraystretch}{1.3} %spécification de l'espace vertical
	\begin{tabular}{|p{3.6cm}|p{6.2cm}|p{6.2cm}|} %largeur des colonnes
		\hline
		\textbf{Type d'approche} & \textbf{Performances appliquées au flux vidéo} & \textbf{Limites} \\
		\hline
		\textbf{Vision par ordinateur classique} 
		& - Détection de changements de plan \newline
		- Calcul d'histogrammes et variations de couleurs \newline
		- Détection basique de mouvement par seuils
		& - Détection strictement syntaxique, sans extraction du sens \newline
		- Absence de reconnaissance d'objets ou de concepts \newline
		- Sensibilité à tous les types de mouvement, y compris les mouvements involontaires \\
		\hline
		\textbf{Apprentissage profond : CNN} 
		& - Classification d'images fixes extraites du flux vidéo \newline
		- Détection d'objets ou caractéristiques dans l'image
		& - Absence de compréhension des relations contextuelles entre les images \newline
		- Incapacité à saisir les dépendances spatio-temporelles \\
		\hline
		\textbf{Apprentissage profond : Transformers} 
		& - Segmentation sémantique et temporelle du flux \newline
		- Analyse du contexte narratif et des relations longues
		& - Nécessite des infrastructures lourdes à la puissance de calcul conéquente \newline
		- Nécessite un entraînement exigeant fondé sur des jeux de données volumineux et adaptés \\
		\hline
	\end{tabular}
	\caption{Comparatif des approches de traitement audiovisuel automatique}
	\label{tab:comparatif_avp}
\end{table}

Comme l'illustre le Tableau~\ref{tab:comparatif_avp}, si les architectures fondées sur les \gls{Transformers} offrent des compétences puissantes, notamment pour l'automatisation de la segmentation temporelle, ces architectures reposent toutefois sur de lourdes puissances de calcul ainsi que sur des jeux de données d'entraînement conséquents. À l'inverse, la \gls{Vision par ordinateur} classique offre une stabilité d'analyse de la syntaxe audiovisuelle, qui ne permet pas toutefois d'en extraire le sens. Ce constat justifie l'orientation de notre étude vers une approche hybride, combinant une architecture stable fondée sur la \gls{Vision par ordinateur} et sur des \gls{cnn} légers, afin de tenter de comprendre dans quelle mesure ces modèles peuvent segmenter un flux vidéo et en extraire des caractéristiques sémantiques.

% TODO: \usepackage{graphicx} required
%\begin{figure}
%	\centering
%	\includegraphics[width=0.7\linewidth]{images/screenshot001}
%	\caption{mahony2020}
%	\label{fig:screenshot001}
%\end{figure}


%transition vers la prochaine sous partie : parler des données d'entrainement (cf fichier IAvideo pattichis)

%Toutes les tâches automatiques qui peuvent être appliquées à la vidéo nécessitent un entraînement multimodal : 
%entrainement multimodal :  When we refer to multimodal or cross-modal learning, we are referring to learning using video, audio and/or text. These approaches generalize to a greater variety of visual tasks including action step localization, action segmentation, text-to-video retrieval, video captioning. \footcite{schiappa2023}
